% !TeX root = relatorio.tex
\documentclass[12pt,a4paper]{article}

\usepackage[T1]{fontenc}
\usepackage[utf8]{inputenc}
\usepackage[brazil]{babel}
\usepackage{lmodern}
\usepackage[a4paper,margin=2.5cm]{geometry}
\usepackage{amsmath}
\usepackage{graphicx}
\usepackage{float}
\usepackage{hyperref}
\usepackage{xcolor}
\usepackage{enumitem}

\hypersetup{
  colorlinks=true,
  linkcolor=blue,
  urlcolor=blue
}

\setlength{\parindent}{0pt}
\setlength{\parskip}{0.6em}

\newcommand{\placeholder}[1]{\textcolor{red}{[Preencher: #1]}}

\begin{document}

\begin{titlepage}
  \centering
  {\large Universidade Estadual de Campinas\par}
  {\large Instituto de Computacao\par}
  \vspace{1.5cm}
  {\Large Introducao ao Processamento Digital de Imagem (MC920/MO443)\par}
  \vspace{0.5cm}
  {\LARGE \textbf{Relatorio do Trabalho 1}\par}
  \vspace{1.5cm}

  \begin{flushleft}
    \textbf{Aluno:} Gustavo Shoiti Sonoda\\
    \textbf{Professor:} Helio Pedrini\\
    \textbf{Data:} \today
  \end{flushleft}

  \vfill
  {\large Campinas -- SP\par}
\end{titlepage}

\tableofcontents
\newpage

\section{Introducao}

Este relatorio documenta a implementacao do Trabalho 1 da disciplina MC920/MO443, cujo objetivo e realizar operacoes basicas de processamento digital de imagens em arquivos PNG. O foco do texto e registrar, para cada item do enunciado, a estrategia adotada, as estruturas de dados utilizadas, os testes executados, os resultados obtidos e as limitacoes observadas.

\placeholder{Escrever um paragrafo curto resumindo a organizacao do programa e a motivacao geral do trabalho.}

\section{Ambiente e Organizacao da Solucao}

\subsection{Ambiente de desenvolvimento}

\begin{itemize}[leftmargin=*]
  \item \textbf{Sistema operacional alvo:} Linux, conforme especificado no enunciado.
  \item \textbf{Formato de entrada e saida:} imagens PNG.
  \item \textbf{Linguagem e bibliotecas:} \placeholder{Informar linguagem escolhida, bibliotecas utilizadas e versoes relevantes.}
  \item \textbf{Compilacao e execucao:} \placeholder{Descrever comandos principais para compilar e executar o projeto.}
\end{itemize}

\subsection{Estruturas de dados e organizacao do codigo}

\begin{itemize}[leftmargin=*]
  \item \textbf{Representacao das imagens:} \placeholder{Explicar como imagens monocromaticas e coloridas sao armazenadas em memoria.}
  \item \textbf{Modularizacao:} \placeholder{Descrever como o codigo foi separado por operacao, modulo ou arquivo.}
  \item \textbf{Convencoes adotadas:} \placeholder{Comentar sobre ordem dos indices, faixas de intensidade, clipping e validacoes.}
  \item \textbf{Leitura e escrita:} \placeholder{Explicar o fluxo de leitura de PNG, processamento e geracao das saidas.}
\end{itemize}

\subsection{Estrategia de testes}

\begin{itemize}[leftmargin=*]
  \item \textbf{Conjunto de imagens:} \placeholder{Listar imagens utilizadas nos experimentos e justificar a escolha.}
  \item \textbf{Criterios de validacao:} \placeholder{Explicar como cada operacao foi conferida visualmente ou numericamente.}
  \item \textbf{Organizacao das saidas:} \placeholder{Indicar onde foram salvos os resultados e como foram nomeados.}
\end{itemize}

\section{Implementacao e Resultados}

\subsection{Item 1.1 -- Rotacao de Imagens em Multiplos de 90 graus}

\begin{itemize}[leftmargin=*]
  \item \textbf{Objetivo:} rotacionar uma imagem monocromatica em 90, 180 e 270 graus sem usar funcoes prontas de rotacao.
  \item \textbf{Implementacao:} \placeholder{Descrever o mapeamento de indices para cada angulo e como foi tratada a dimensao da imagem.}
  \item \textbf{Testes executados:} \placeholder{Comparar imagem original com as tres rotacoes e comentar como a corretude foi verificada.}
  \item \textbf{Resultados e discussao:} \placeholder{Inserir observacoes sobre custo, simplicidade da solucao e possiveis erros comuns evitados.}
\end{itemize}

\subsection{Item 1.2 -- Ampliacao de Imagem por Replicacao}

\begin{itemize}[leftmargin=*]
  \item \textbf{Objetivo:} ampliar uma imagem monocromatica por fatores 2 e 4 usando replicacao de pixels.
  \item \textbf{Implementacao:} \placeholder{Explicar como cada pixel original gera um bloco na imagem ampliada.}
  \item \textbf{Testes executados:} \placeholder{Descrever os casos usados para conferir dimensoes e preservacao do padrao visual.}
  \item \textbf{Resultados e discussao:} \placeholder{Comentar o efeito serrilhado esperado e a diferenca em relacao a interpolacao.}
\end{itemize}

\subsection{Item 1.3 -- Esboco a Lapis}

\begin{itemize}[leftmargin=*]
  \item \textbf{Objetivo:} produzir um efeito de esboco a lapis a partir de imagem colorida.
  \item \textbf{Implementacao:} \placeholder{Descrever as etapas de conversao para cinza, desfoque gaussiano e divisao entre imagens.}
  \item \textbf{Parametros adotados:} mascara gaussiana de referencia com dimensao $21 \times 21$ pixels.
  \item \textbf{Testes executados:} \placeholder{Explicar quais imagens evidenciaram melhor contornos, sombras e texturas.}
  \item \textbf{Resultados e discussao:} \placeholder{Comentar o efeito visual obtido e a influencia do desfoque no destaque de bordas.}
\end{itemize}

\subsection{Item 1.4 -- Ajuste de Brilho}

\begin{itemize}[leftmargin=*]
  \item \textbf{Objetivo:} ajustar o brilho de uma imagem monocromatica por correcao gama.
  \item \textbf{Implementacao:} \placeholder{Descrever a normalizacao para $[0,1]$, a aplicacao da equacao de gama e o retorno para $[0,255]$.}
  \item \textbf{Formula utilizada:} \placeholder{Registrar explicitamente a equacao implementada, por exemplo $B = A^{1/\gamma}$, caso esse tenha sido o formato usado no codigo.}
  \item \textbf{Testes executados:} \placeholder{Listar os valores de $\gamma$ usados, por exemplo 1.5, 2.5 e 3.5, e comparar os efeitos.}
  \item \textbf{Resultados e discussao:} \placeholder{Comentar como diferentes valores de gama escurecem ou clareiam a imagem.}
\end{itemize}

\subsection{Item 1.5 -- Binarizacao por Limiar}

\begin{itemize}[leftmargin=*]
  \item \textbf{Objetivo:} gerar uma imagem binaria a partir de uma imagem monocromatica e de um limiar fixo $T$.
  \item \textbf{Implementacao:} \placeholder{Descrever a regra de conversao para 0 e 255 e como o limiar foi parametrizado.}
  \item \textbf{Formula utilizada:} $B = 255$ se $A > T$; caso contrario, $B = 0$.
  \item \textbf{Testes executados:} \placeholder{Informar os limiares avaliados e justificar a escolha do valor final apresentado.}
  \item \textbf{Resultados e discussao:} \placeholder{Discutir perdas de detalhe, ruido e separacao entre objeto e fundo.}
\end{itemize}

\subsection{Item 1.6 -- Mosaico}

\begin{itemize}[leftmargin=*]
  \item \textbf{Objetivo:} construir um mosaico $4 \times 4$ a partir de uma imagem monocromatica, seguindo a numeracao indicada no enunciado.
  \item \textbf{Implementacao:} \placeholder{Explicar como a imagem foi dividida em blocos e como a nova ordem foi aplicada.}
  \item \textbf{Testes executados:} \placeholder{Descrever como foi conferida a correspondencia entre a ordem original dos blocos e a ordem remapeada.}
  \item \textbf{Resultados e discussao:} \placeholder{Comentar sobre o efeito visual do embaralhamento e sobre cuidados com divisao exata em blocos.}
\end{itemize}

\subsection{Item 1.7 -- Alteracao de Cores}

\begin{itemize}[leftmargin=*]
  \item \textbf{Objetivo:} simular o efeito de fotografias antigas por meio de uma transformacao linear nas componentes RGB.
  \item \textbf{Implementacao:} \placeholder{Descrever a multiplicacao de cada pixel pela matriz de transformacao e o clipping para o intervalo $[0,255]$.}
  \item \textbf{Formula utilizada:} \placeholder{Inserir a matriz utilizada no relatorio final, conforme o enunciado.}
  \item \textbf{Testes executados:} \placeholder{Indicar imagens usadas para avaliar tons de pele, vegetacao e contraste.}
  \item \textbf{Resultados e discussao:} \placeholder{Comentar o efeito visual do filtro e o impacto da saturacao em 255.}
\end{itemize}

\subsection{Item 1.8 -- Transformacao de Imagens Coloridas}

\begin{itemize}[leftmargin=*]
  \item \textbf{Objetivo:} realizar operacoes adicionais sobre imagens RGB, incluindo (a) a transformacao componente a componente e (b) a geracao de uma unica banda pela media ponderada.
  \item \textbf{Implementacao do item (a):} \placeholder{Descrever a aplicacao direta das equacoes para $R'$, $G'$ e $B'$ e o tratamento de overflow.}
  \item \textbf{Implementacao do item (b):} usar a intensidade $I = 0.2989R + 0.5870G + 0.1140B$ para obter uma imagem com uma unica banda.
  \item \textbf{Testes executados:} \placeholder{Explicar como foi validada a saida colorida e a conversao para intensidade unica.}
  \item \textbf{Resultados e discussao:} \placeholder{Comparar a saida colorida transformada com a versao em tons de cinza obtida pela media ponderada.}
\end{itemize}

\subsection{Item 1.9 -- Planos de Bits}

\begin{itemize}[leftmargin=*]
  \item \textbf{Objetivo:} extrair os planos de bits de uma imagem monocromatica.
  \item \textbf{Implementacao:} \placeholder{Explicar como cada pixel foi decomposto em base 2 e como os planos individuais foram reconstruidos.}
  \item \textbf{Formula utilizada:} \placeholder{Registrar a representacao polinomial em base 2 apresentada no enunciado.}
  \item \textbf{Testes executados:} \placeholder{Descrever a analise de planos como bit 0, bit 4 e bit 7.}
  \item \textbf{Resultados e discussao:} \placeholder{Comentar a contribuicao visual de cada plano para a imagem final.}
\end{itemize}

\subsection{Item 1.10 -- Combinacao de Imagens}

\begin{itemize}[leftmargin=*]
  \item \textbf{Objetivo:} combinar duas imagens monocromaticas de mesmo tamanho por media ponderada.
  \item \textbf{Implementacao:} \placeholder{Inserir a expressao usada para combinar as imagens e explicar a escolha dos pesos.}
  \item \textbf{Testes executados:} \placeholder{Descrever o conjunto de pares de imagens usados e o que foi observado para pesos extremos e intermediarios.}
  \item \textbf{Resultados e discussao:} \placeholder{Comentar mistura visual, suavizacao e preservacao de detalhes.}
\end{itemize}

\subsection{Item 1.11 -- Transformacao de Intensidade}

\begin{itemize}[leftmargin=*]
  \item \textbf{Objetivo:} aplicar diferentes transformacoes sobre imagens monocromaticas, incluindo reducao do numero de niveis de cinza, negativo, remapeamento do intervalo de intensidades, reversao das linhas pares, espelhamento da metade superior na metade inferior e espelhamento vertical completo.
  \item \textbf{Implementacao:} \placeholder{Descrever cada subitem separadamente, incluindo a quantizacao para diferentes numeros de niveis como 256, 64, 32, 16, 8, 4 e 2.}
  \item \textbf{Testes executados:} \placeholder{Explicar como cada transformacao foi validada visualmente.}
  \item \textbf{Resultados e discussao:} \placeholder{Comparar as perdas de detalhe da quantizacao e os efeitos geometricos das operacoes de espelhamento e reversao.}
\end{itemize}

\subsection{Filtros espaciais adicionais}

\begin{itemize}[leftmargin=*]
  \item \textbf{Contexto:} o enunciado apresenta mascaras $h_1$ ate $h_{11}$ e solicita discutir os efeitos de cada filtro.
  \item \textbf{Implementacao:} \placeholder{Explicar como foi feita a convolucao, o tratamento de borda e a normalizacao das saidas.}
  \item \textbf{Combinacao de filtros:} os filtros $h_3$ e $h_4$ devem ser aplicados individualmente e tambem de forma combinada por meio de $\sqrt{(h_3)^2 + (h_4)^2}$.
  \item \textbf{Resultados e discussao:} \placeholder{Descrever, para cada mascara, se o efeito principal foi suavizacao, realce, deteccao de bordas ou outro comportamento observado.}
\end{itemize}

\section{Discussao Geral}

\begin{itemize}[leftmargin=*]
  \item \textbf{Comparacao entre operacoes pontuais e espaciais:} \placeholder{Discutir diferencas entre operacoes por pixel, reordenacoes geometricas e filtragem por vizinhanca.}
  \item \textbf{Desempenho:} \placeholder{Comentar tempo de execucao, gargalos e oportunidades de vetorizacao.}
  \item \textbf{Qualidade visual dos resultados:} \placeholder{Resumir quais operacoes tiveram melhor resultado e em quais tipos de imagem.}
\end{itemize}

\section{Limitacoes}

\begin{itemize}[leftmargin=*]
  \item \placeholder{Listar restricoes da implementacao, como tratamento de borda, parametros fixos, dependencias de tamanho ou sensibilidade a ruido.}
  \item \placeholder{Indicar situacoes especiais nao tratadas pelo programa, conforme solicitado no enunciado.}
\end{itemize}

\section{Conclusao}

Neste trabalho foram implementadas diversas operacoes de processamento digital de imagens envolvendo transformacoes geometricas, operacoes pontuais, manipulacao de cor, extracao de informacao binaria e filtragem espacial. O conjunto de resultados apresentado ao longo do relatorio permite discutir tanto a corretude funcional da implementacao quanto os efeitos visuais produzidos por cada tecnica.

\placeholder{Escrever um paragrafo final sintetizando os principais aprendizados, os itens mais desafiadores e os resultados mais relevantes.}

\end{document}
